\label{chap:evaluation}
This chapter answers Research Question~3:

\begin{quote}
\textit{How do students plan with the hybrid interface, and what does the graph-based view contribute to their planning?}
\end{quote}

It reports the evaluation study: the study's goal and the questions that guided it (Section~\ref{sec:eval-design}), the system as the artefact under evaluation (Section~\ref{sec:eval-system}), the participants (Section~\ref{sec:eval-participants}), the procedure (Section~\ref{sec:eval-procedure}) and the analysis method (Section~\ref{sec:eval-analysis}), then the findings (Section~\ref{sec:eval-results}), which Chapter~\ref{chap:discussion} interprets. In the terms of the nested model, the study validates the encoding, interaction and algorithm levels of the artefact~\cite{munzner2009}.

\section{Study Goal and Guiding Questions}
\label{sec:eval-design}

The study investigates how a curriculum graph alters the planning process of students equipped with an established tabular workflow. While the formative study treated the graph view's utility as a premise (Section~\ref{sec:overall-assumptions}) and conceptually partitioned the interface into exploratory (graph) and commitment-oriented (table) modalities (Section~\ref{sec:mutual-exclusion}), neither assumption had undergone evaluation. We therefore conducted a within-subjects observational study ($N=11$) where participants completed comparable planning tasks using the system of Chapter~\ref{chap:implementation} first with the tabular baseline alone, then with the graph view enabled. Granular, frame-by-frame behavioral logging paired with participant reports allows the graph's operational contribution to be traced directly within the planning process. This focus on process tracing governed the research design. Detecting reliable outcome differences in speed or plan quality would require a high-powered, counterbalanced sample matched for domain knowledge; we deliberately prioritized observational fidelity over broad generalization. Across 22 recorded sessions, frame-level coding generated descriptive findings and grounded hypotheses. We used a fixed condition order intentionally within this framing; its methodological trade-offs were anticipated (Section~\ref{sec:eval-order}), and no empirical claims in this chapter depend on comparative inferences that could be compromised by order effects.

\subsection{Guiding Questions}
\label{sec:eval-guiding}

To operationalise this goal, the study was guided by four questions, carried over from the proposal for this thesis, which planned a task-based evaluation of the hybrid interface and of the recommendations embedded in it:

\begin{itemize}
    \item \textbf{Planning practice:} How do students execute the planning process using the tool?
    \item \textbf{The graph's contribution:} What does the graph-based view add to that work? Answered by comparing the two conditions within each participant (Section~\ref{sec:eval-ab}).
    \item \textbf{Recommendations in the workflow:} Where and how do the embedded recommendations help?
    \item \textbf{Outcomes and satisfaction:} Do the plans students produce meet the task they were set, and how satisfied are students with the tool?
\end{itemize}

\subsection{Scope and Non-Goals}
\label{sec:eval-nongoals}

This study does not evaluate plan accuracy or generation speed, nor does it use either to compare the two conditions. Establishing reliable outcome measures would require matched cohorts with a known baseline of curriculum knowledge, a tool-agnostic operationalization of plan quality, and a sample large enough to isolate condition effects from individual variance. These prerequisites fall outside the scope of a qualitative, small-scale observational design. The completion audit in Section~\ref{sec:eval-effectiveness} is therefore not a performance comparison. It serves a single methodological purpose: by independently establishing which plans met the brief, it reveals the gap between perceived and actual success, thereby contextualizing the self-reported data.

\section{System as Evaluation Artefact}
\label{sec:eval-system}

The evaluated system is the tool detailed in Chapters~\ref{chap:design} and~\ref{chap:implementation}. To provide context for the analysis codes, the components carrying the findings are: the Table View (Section~\ref{sec:table-view}), the Graph View with its filter sidebar and prerequisite overlay (Section~\ref{sec:graph-spatial}), the Dashboard (Sections~\ref{sec:dashboard} and~\ref{sec:compliance}), the Recommendation Panel (Section~\ref{sec:rec-panel}), the Parking Stage (Section~\ref{sec:parking}), the Course Catalogue (Section~\ref{sec:sidebar-design}), the visual legend (Section~\ref{sec:legend}), and the onboarding tour (Section~\ref{sec:onboarding}).

\section{Participants and Recruitment}
\label{sec:eval-participants}
%% REV (is this kind of table with the participants missing in chapter 4?)
Twelve sessions were conducted and eleven are analysed. One session (P08) was excluded due to a lack of consent for the screen recording, which precluded the visual analysis method detailed in Section~\ref{sec:eval-analysis}.

Participants were recruited from the Faculty of Informatics at TU~Wien through the author's network. Eligibility was restricted to students enrolled in one of the two target programmes; no compensation was provided. The eleven analysed participants ranged from their 2nd to 8th semester across Software Engineering (MSc) and Computer Science (BSc), and they span the target experience strata: four reported regular prior experience with study-planning tools, one reported little, and six reported none (Table~\ref{tab:participants}). No recording frames are reproduced in this thesis, and participants are referred to exclusively by assigned codes.

\begin{table}[htpb]
    \centering
    \small
    \begin{tabular}{llcl}
        \toprule
        P & Programme & Sem. & Prior exp. \\
        \midrule
        P01 & Software Engineering (MSc) & 6 & none    \\
        P02 & Software Engineering (MSc) & 6 & regular \\
        P03 & Computer Science (BSc)     & 6 & none    \\
        P04 & Software Engineering (MSc) & 5 & regular \\
        P05 & Software Engineering (MSc) & 8 & none    \\
        P06 & Software Engineering (MSc) & 2 & none    \\
        P07 & Software Engineering (MSc) & 8 & regular \\
        P09 & Software Engineering (MSc) & 5 & little  \\
        P10 & Computer Science (BSc)     & 2 & none    \\
        P11 & Computer Science (BSc)     & 2 & none    \\
        P12 & Software Engineering (MSc) & 6 & regular \\
        \bottomrule
    \end{tabular}
    \caption{The eleven analysed participants. P08 is excluded and not listed.}
    \label{tab:participants}
\end{table}

\section{Procedure}
\label{sec:eval-procedure}

Each participant used the tool under two conditions in a fixed order during a single session. This section details the conditions and their task briefs, the rationale for the fixed order and its methodological trade-offs, and the session sequence alongside its measurement instruments.

\subsection{Conditions and Scenarios}
\label{sec:eval-conditions}

Each participant completed two planning tasks under different graph availability settings. To control for individual domain preferences, these tasks were framed as persona scenarios. The graph setting was an account-level property fixed before each task began; participants could not toggle the view on or off during a scenario.

\begin{itemize}
    \item \textbf{Scenario~A (graph off).} A Cybersecurity persona, planned with the graph unavailable. The table, Dashboard, recommendations, catalogue, and Parking Stage remain fully available.
    \item \textbf{Scenario~B (graph on).} An Artificial Intelligence persona, planned with the full tool including the graph and its filter sidebar.
\end{itemize}

Everything other than graph availability was held constant across the two conditions. Each brief stated only its persona constraints, namely the per-semester ECTS band and the single reduced-load semester, along with strategic guidance on preferred module types. The curriculum's inherent hard requirements (the mandatory modules, the chosen focus area, the narrow (\emph{enge Wahl}) and broad (\emph{breite Wahl}) electives, and the introductory-phase (StEOP) rules) were intentionally excluded from the brief. Participants encountered these only inside the tool and tracked them via the Dashboard, making the Dashboard the sole channel through which curriculum constraints were communicated during a task. The two briefs are reproduced exactly as presented in %% REV (this is not a section isnt it?) 
Section~\ref{app:eval-scenarios}.

\subsection{Fixed Scenario Order}
\label{sec:eval-order}

Scenario~A preceded Scenario~B for every participant. We opted for a fixed sequence rather than counterbalancing because the object of study is how a graph enters a workflow a student has already formed. By fixing the order, every participant encountered the graph at the same point on their learning curve. While counterbalancing is the standard approach to isolate treatment effects, it would have yielded two different baseline workflows within the cohort (one formed with the graph, one without). We chose to prioritize a consistent baseline workflow, grounding this decision in Greenwald's analysis: %% REV (please check in detail if this source can ground this claim because this is crucual here)
counterbalancing distributes practice effects but does not remove an interaction between treatment and practice~\cite[p.~320]{greenwald1976within}, and a fixed sequence is an appropriate alternative when the course of learning is itself the object of study~\cite[pp.~316--318]{greenwald1976within}.

Because graph availability and second-task exposure are perfectly collinear, this chapter attributes no changes in speed or outcome to the graph. Instead, Section~\ref{sec:eval-ab} focuses strictly on whether the graph displaces the established workflow, explicitly labelling any inseparable second-run effects.

\subsection{Session Sequence and Instruments}
\label{sec:eval-sequence}

Sessions were screen- and audio-recorded and administered via a custom web application (Section~\ref{sec:impl-userstudy}) following a fixed sequence: (1)~demographics, (2)~consent, (3)~video tutorial, (4)~Scenario~A, (5)~post-task questionnaire, (6)~Scenario~B, (7)~post-task questionnaire, (8)~open exploration, (9)~final questionnaire, and (10)~a semi-structured interview. Participants used the think-aloud method during all working phases, with a moderator present to resolve operational questions.

The post-task questionnaires contained nine items: three assessing task success, plan satisfaction, and difficulty, alongside six domain-adapted items from the Technology Acceptance Model (TAM) measuring perceived usefulness and ease of use. The final questionnaire consisted of the standard 26-item User Experience Questionnaire (UEQ)~\cite{laugwitz_construction_2008}. All instruments are reproduced in Section~\ref{app:eval-instruments}.
\TDrev{How the TAM-derived items are reported}{Moritz, 2026-09-06: to be discussed. The items are a three-per-construct, reworded short form, so Davis's reliabilities cannot be claimed. Options: keep reporting them as TAM-derived means without a reliability claim, as now; compute Cronbach's alpha over the eleven responses per condition and report it with the caveat that eleven cases support only a rough estimate; or drop the construct labels and report the six items descriptively. The choice also touches Chapter~\ref{chap:discussion}'s limitation on the instrument.}
%% REV (look how i changed it, so i think we are saying we used 6 items form tam and adopted it to the domain?)

\section{Analysis Method}
\label{sec:eval-analysis}

This section details the procedures used to process and analyze the study data, encompassing the preparation of dual data records, the rationale for fixed-interval frame sampling, and the application of a structured thematic codebook.

\subsection{Two Records per Participant}
\label{sec:eval-tworecords}

Two independent data streams were collected and triangulated per participant (Figure~\ref{fig:eval-pipeline}): screen recordings and audio transcripts. The audio recordings were transcribed using %% REV Please make the same footnote like in chapter 4)
Otter.ai Screen recordings were sampled at ten-second intervals, yielding 2{,}179 frames across the 22 scenario runs. Each frame underwent descriptive coding against a fixed scheme capturing the active surface, open panels, sidebar contents, per-semester ECTS, Dashboard state, and transient banners. Frame coding was executed by a locally hosted vision-language model (\emph{Qwen3-VL-30B-A3B-Instruct}), ensuring data privacy. The model strictly adhered to the protocol detailed in Section~\ref{app:eval-stagelog}: recording one independent row per frame without interpolation, marking occluded values as unreadable, and cross-referencing all plan-state metrics against the Dashboard. To validate the model's reliability, the author manually coded three complete runs; agreement metrics are reported in Section~\ref{app:eval-agreement}. Crucially, all interpretative analysis, severity assessment, and completion auditing were conducted exclusively by the author.

\subsection{Rationale and Limitations of Frame Sampling}
\label{sec:eval-framemethod}

Systematic fixed-interval frame sampling was chosen to ensure objective, gapless comparability of interface occupancy and dwell times across the cohort. Because the label scheme was fixed prior to analysis, the process states of Section~\ref{sec:eval-statemodel} are empirically derived from the data rather than selectively interpreted by the analyst. %% REV (add that we also did this because limmited time constrains within this study, handycoding all states for these videos would take whay longer)

However, this quantitative lens has inherent methodological trade-offs. First, a ten-second resolution inevitably misses sub-interval transient events (such as brief hovers or quickly dismissed menus), making these visual counts lower bounds. Second, a visual frame captures interface state but not user intent. To mitigate this, the visual codes are triangulated with the transcripts, which potentially carry the underlying rationale for observed actions. Third, complex UI configurations (such as simultaneous open panels) required uniform compound-state rules (Table~\ref{tab:states}) which dictate the resulting occupancies. Finally, the method's strict reliance on visual data necessitated the exclusion of P08, for whom no screen recording was available. These limitations define their interpretative boundaries as carried forward in Chapter~\ref{chap:discussion}.

\subsection{Codebook and Frequencies}
\label{sec:eval-codebook}

Transcripts and recordings were analyzed using codebook thematic analysis~\cite{braun_thematic_2022}. Unlike the formative study's inductive approach, this evaluation utilized a deductive, controlled codebook to classify usability problems (\mbox{E-P}), positive findings (\mbox{E-G}), and neutral or methodological observations (\mbox{E-N}). Each code maps to an %% REV (reference?)
ISO~9241-11 dimension, and usability problems were assigned a severity rating on Nielsen's five-point scale~\cite{nielsen1993}. To ensure cohort-wide comparability, any new code identified during later sessions was retroactively evaluated for all prior sessions. The resulting matrix is reported in Section~\ref{app:eval-matrix}.

Frequencies are reported in the format \emph{n}/\emph{m}, where \emph{m} is the number of participants for whom a code could be definitively evaluated (e.g., they reached the relevant feature), and \emph{n} is the number of participants for whom the code was evidenced. Thus, a code reported as 9/9 indicates the behavior occurred in all nine sessions where the context arose, not in nine out of eleven total sessions. Because severity is a single-analyst judgment, the ranking in Table~\ref{tab:problems} is indicative rather than exact.

\begin{figure}[htbp]
\centering
\begin{tikzpicture}[font=\small,
  b/.style={draw, line width=1pt, rounded corners=3pt, fill=white, align=center, inner sep=2.4mm, text width=3.9cm},
  wide/.style={b, text width=9.0cm},
  a/.style={-{Latex[length=2.8mm]}, line width=1.1pt, draw=black}]
  \node[b] (rec) at (-2.9,0) {\textbf{Screen recording}};
  \node[b] (tra) at ( 2.9,0) {\textbf{Transcript} \\ \small think-aloud, interview};
  \node[b] (frm) at (-2.9,-1.9) {\textbf{Frames} \\ \small every 10\,s};
  \node[b] (qual) at ( 2.9,-1.9) {\textbf{Utterance coding}};
  \node[b] (log) at (-2.9,-3.8) {\textbf{State log} \\ \small surface, panels, \\ \small lane ECTS};
  \node[wide] (cb) at (0,-5.9) {\textbf{Codebook applied to both records, for every participant}};
  \node[wide] (out) at (0,-7.7) {\textbf{Process model, code frequencies, completion audit}};
  \draw[a] (rec) -- (frm);
  \draw[a] (tra) -- (qual);
  \draw[a] (frm) -- (log);
  % Each record drops straight into its own side of the codebook box, so the two
  % paths never cross and neither arrowhead lands on the other's line.
  \draw[a] (log.south)  -- ([xshift=-2.9cm]cb.north);
  \draw[a] (qual.south) -- ([xshift=2.9cm]cb.north);
  \draw[a] (cb) -- (out);
\end{tikzpicture}
\caption{The two-record analysis. The screen recording is sampled and labelled independently of the transcript, and the codebook is then applied to both records for every participant.}
\label{fig:eval-pipeline}
\end{figure}

% =============================================================================

\section{Findings}
\label{sec:eval-results}

The findings are reported in the order of the guiding questions: the planning process first, then what the graph changes in it, then the surfaces around the loop, including the Recommendation Panel, and what worked and where the loop breaks, then the questionnaires and the plans produced against self-report, and a summary that Chapter~\ref{chap:discussion} takes up. %% REV (Frequencies?) 
Frequencies are written as in Section~\ref{sec:eval-codebook}.

\subsection{The Planning Loop}
\label{sec:eval-process}

Planning under this tool is a single checklist-driven loop, and it has the same shape for every participant. Students read what the Dashboard says is still missing, find a candidate that closes the gap, place it, and return to the list. This section establishes that loop from the state logs.

\paragraph{A State Model of the Planning Process}
\label{sec:eval-statemodel}
%% REV (in limitation should be clear that we cannot caption indent: "opnen all the tome but never use it") 
%% REV (how could we argue that the recommendation panel is not a state here?)
Each sampled frame carries a primary surface, and, recorded independently of it, which
panels were open: the Dashboard, which coexists with either view rather than replacing
it; the sidebar, which holds the Course Catalogue in the Table View and the filter panel
in the Graph View, and can be hidden in the table; the Recommendation Panel; the visual
legend; and any modal. The process states of Table~\ref{tab:states} combine the surface
with the two that differentiate it, the Dashboard and the sidebar. They differentiate only
within the Table View: the Graph View's sidebar is always the filter panel, so it yields
a single state.

\begin{table}[htpb]
    \centering
    \small
    \begin{tabular}{lp{7.4cm}}
        \toprule
        State & Interpretation \\
        \midrule
        \textsc{dash+cat} & Table with the Dashboard open and the catalogue in the sidebar: the checklist-driven working state \\
        \textsc{cat}      & Table with the catalogue, Dashboard closed: browsing for candidates \\
        \textsc{dash}     & Table with the Dashboard, no catalogue: reading status \\
        \textsc{graph}    & Graph View; its sidebar is always the filter panel \\
        \textsc{table}    & Table with neither panel \\
        \textsc{off}      & Study guide, sign-in, loading, questionnaire \\
        \bottomrule
    \end{tabular}
    \caption{Process states derived from the frame labels. The compound definition is required because the Dashboard is a right-hand panel that coexists with the table rather than replacing it.}
    \label{tab:states}
\end{table}

%% REV (this table still shpows the rec overlay, but doest talk about it.)
\begin{table}[htpb]
    \centering
    \small
    \begin{tabular}{lrrr@{\hskip 2em}rrr}
        \toprule
        & \multicolumn{3}{c}{Scenario~A (1{,}197 frames)} & \multicolumn{3}{c}{Scenario~B (982 frames)} \\
        \cmidrule(r){2-4}\cmidrule(l){5-7}
        State & Occ.\ \% & Entries & Dwell~s & Occ.\ \% & Entries & Dwell~s \\
        \midrule
        \textsc{dash+cat} & 51.2 & 38 & 161 & 46.5 & 51 & 90 \\
        \textsc{cat}      & 27.1 & 40 &  81 &  6.7 & 22 & 30 \\
        \textsc{graph}    & ---  & ---& --- & 26.6 & 42 & 62 \\
        \textsc{dash}     &  5.5 & 11 &  60 &  8.2 &  8 & 101 \\
        \textsc{table}    &  0.6 &  5 &  14 &  2.0 &  3 & 67 \\
        \textsc{off}      & 15.6 & 26 &  72 &  9.9 & 35 & 28 \\
        \midrule
        \multicolumn{7}{l}{\emph{Overlay, counted independently of the states above}} \\
        \textsc{recs}     & 24   & 23 & 125 & 20   &  8 & 244 \\
        \bottomrule
    \end{tabular}
    \caption{State occupancy, entries and mean dwell per entry over the eleven runs of each condition. The catalogue and the graph exchange places: 40 catalogue entries in Scenario~A against 42 graph entries in Scenario~B, at comparable dwell. The Recommendation Panel is an overlay; its occupancy is a share of the same frames and is outside the column total.}
    \label{tab:occupancy}
\end{table}

\paragraph{The Core Workflow: A Checklist-Driven Two-State Loop}
\label{sec:eval-loop}

The Dashboard's missing-requirements checklist was the primary driver of task completion for every participant (E-G15, 11/11). Students consistently seeded their plans with the pre-filled mandatory backbone, a feature universally accepted and valued (E-G01, 11/11). From this baseline, they utilized the checklist to identify remaining gaps, took action to close them, and returned to the list to verify progress. For instance, P09 began Scenario~A by reading the count aloud, \enquote{22 missing requirements, wow, that is a lot}, and then systematically resolved them item by item. Similarly, P10 explicitly articulated this feedback loop: \enquote{I need narrow-elective modules, transferable skills, so let me see what narrow-elective ones there are}. For the most junior participant, this loop only emerged after moderator intervention: P06 had been placing courses for roughly nine minutes with no view of what remained, and his response upon seeing the tracker marked the transition from unstructured browsing to goal-directed planning (\enquote{exactly that, here I see what I need}). That this loop takes the exact same shape for independent and heavily assisted participants demonstrates that it is a structural property of the tool rather than an artefact of moderation.

The loop alternates between two states: the checklist-reading state (\textsc{dash+cat})
and a candidate-finding state, which in Scenario~A is the catalogue. The transition
\textsc{cat}~$\rightarrow$~\textsc{dash+cat} accounts for 27.5\,\% of all state changes in
Scenario~A and the reverse for 17.4\,\%, so this single pair constitutes 44.9\,\% of all
logged state transitions. The cycle is tight: the median path length from an open
Dashboard back to the Dashboard is 2.0 transitions ($n=38$ returns), the shortest a round
trip can be once self-transitions are excluded. Students read the list, made one excursion
to find candidates, and returned to the list.

\subsection{What the Graph Changes}
\label{sec:eval-ab}

The graph alters the established planning loop at exactly one point: it replaces the catalogue as the primary interface for translating a missing requirement into candidate courses, leaving all downstream actions unchanged. This section details this workflow substitution, the resulting division of labour between the two views, and the system rejection rate, which remained constant. Finally, practice effects confounded by the fixed-order design are isolated at the end of the section and are not attributed to the graph.

\paragraph{The Graph Replaces the Catalogue}
%% REV (it should be clear that there was no force to use the graph, they got told you have it now, thats it)
\label{sec:eval-substitution}
Rather than adding a third station to the planning cycle, the graph displaces the catalogue entirely. In Scenario~B, the feedback loop shifts to alternate between the Dashboard and the graph. The \textsc{dash+cat}~$\rightarrow$~\textsc{graph} transition and its reverse each account for 19.3\,\% of all state changes (38.6\,\% combined), against the 44.9\,\% share the catalogue--Dashboard cycle held in Scenario~A. Frame occupancy shows the same substitution: catalogue-only occupancy drops from 27.1\,\% to 6.7\,\%, while graph occupancy rises to 26.6\,\%, close to the share the catalogue vacates. Entry counts are similarly aligned: the catalogue was entered 40 times across the eleven Scenario~A runs, whereas the graph was entered 42 times in Scenario~B, with comparable mean dwell times (81~s versus 62~s). 

Crucially, per sampled frame spent in the checklist-reading state (\textsc{dash+cat}), the chance of moving directly to the catalogue drops from 3.1\,\% in Scenario~A to 0.0\,\% in Scenario~B, against 6.4\,\% for the graph. The catalogue does not disappear in Scenario~B, but the direct path from the checklist is severed: the only route to it from within the planner is the side branch out of the graph, taken six times and drawn at 4.1\,\% of transitions. At the participant level, catalogue occupancy falls in ten out of eleven cases, dropping from a median of 27.4\,\% to 4.5\,\% (Table~\ref{tab:ab}).

Despite this substitution, the cycle maintains its structural integrity. In both conditions, the dominant excursion from the checklist is a direct round trip to a single candidate-finding surface. These reciprocal transitions account for 44.9\,\% of all state changes in Scenario~A and 38.6\,\% in Scenario~B. Furthermore, the median path length from an open Dashboard back to an open Dashboard remains exactly two state changes in both scenarios, the minimum for a two-station cycle. The introduction of the third surface did not lengthen the excursion from the checklist; it redirected it.

Two things separate this shift from a practice effect. The catalogue remained available and unchanged in Scenario~B, and nothing in the task or the instrument asked participants to stop using it; the study guide asked them to use the whole tool including the graph (Section~\ref{sec:eval-conditions}), which accounts for the graph being opened but not for its taking the catalogue's place in the loop rather than joining it. Growing familiarity likewise accounts for a faster second run, not for the abandonment of a surface participants had used competently in the first. The displacement is therefore the comparison in this chapter least exposed to the fixed order, and the extent to which the instruction shapes it is carried forward to Section~\ref{sec:disc-limitations}.

%% REV (lets discuss why the table numbers are so low)
\begin{figure}[htbp]
\centering
\begin{tikzpicture}[font=\small,
  s/.style={draw, line width=1pt, rounded corners=3pt, fill=white, align=center,
            inner sep=2mm, text width=2.9cm, minimum height=12mm},
  faint/.style={s, dashed, draw=gray, text=gray},
  a/.style={-{Latex[length=2.6mm]}, line width=1.1pt, draw=black},
  lb/.style={font=\scriptsize, inner sep=1.4pt, fill=white}]
  % ---- Scenario A: the cycle runs between the catalogue and the Dashboard
  \node[font=\small\bfseries] at (0,1.5) {Scenario A (graph off)};
  \node[s] (dashA) at (0,0)    {\textbf{\textsc{dash+cat}} \\ \small 51.2\,\%};
  \node[s] (catA)  at (0,-3.0) {\textbf{\textsc{cat}} \\ \small 27.1\,\%};
  \draw[a] ([xshift=-0.85cm]dashA.south) -- node[lb, left]  {17.4\,\%} ([xshift=-0.85cm]catA.north);
  \draw[a] ([xshift= 0.85cm]catA.north)  -- node[lb, right] {27.5\,\%} ([xshift= 0.85cm]dashA.south);
  % ---- Scenario B: the same cycle, with the graph in the catalogue's place
  \node[font=\small\bfseries] at (6.6,1.5) {Scenario B (graph on)};
  \node[s] (dashB) at (6.6,0)    {\textbf{\textsc{dash+cat}} \\ \small 46.5\,\%};
  \node[s] (grB)   at (6.6,-3.0) {\textbf{\textsc{graph}} \\ \small 26.6\,\%};
  \draw[a] ([xshift=-0.85cm]dashB.south) -- node[lb, left]  {19.3\,\%} ([xshift=-0.85cm]grB.north);
  \draw[a] ([xshift= 0.85cm]grB.north)   -- node[lb, right] {19.3\,\%} ([xshift= 0.85cm]dashB.south);
  % The catalogue is no longer a station of the cycle: it hangs off the graph as a
  % side branch taken on 4.1 per cent of transitions.
  \node[faint] (catB) at (6.6,-5.9) {\textsc{cat} \\ \small 6.7\,\%};
  \draw[a, draw=gray, dashed] (grB.south) -- node[lb, right, text=gray] {4.1\,\%} (catB.north);
\end{tikzpicture}
\caption{The dominant planning cycle in each condition, with each transition's share of all state changes and each state's occupancy. Shape and period are unchanged; only the candidate-finding surface differs.}
\label{fig:eval-loop}
\end{figure}

\paragraph{Query Support Without Commit Support}
\label{sec:eval-graphrole}

The recordings reveal exactly why the graph displaces the catalogue but does not absorb the entire planning loop. It effectively supports the \emph{query} step (identifying candidate courses for a requirement) but structurally inhibits the \emph{commit} step (assigning a course to a specific semester). This confirms the intended division of labour discussed in Section~\ref{sec:mutual-exclusion}: exploration belongs to the graph, and commitment to the table.
The graph excels during the query phase. For example, P02 read a missing requirement from the Dashboard, switched to the graph, applied the narrow-elective filter, harvested candidates into the Parking Stage, and returned to the table to place them. Others utilized combined filters for advanced queries (E-G18): P09 combined term-availability and \enquote{not planned} filters to deduce a definitive negative answer (\enquote{everything offered in winter and summer would already be planned}), while P12 used the same combination to locate and place twelve credits. Furthermore, P09 explicitly relied on the graph to isolate specific credit values: \enquote{for that I switch to the graph view, I see it immediately with my filters}.

Conversely, the commit step forces a return to the table due to an intentional design constraint. While the graph allows direct course placement, it deliberately omits the semester dimension; neither temporal lanes nor per-semester ECTS totals are rendered. Because users cannot see which semester is being loaded, committing a course directly from a graph node relies on blind working memory. P10 articulated this explicitly (\enquote{adding directly from here I would probably not do, because I do not know which semester has room}), only succeeding later by recalling previous context (\enquote{I could still recall that I wanted to do something there anyway}). P12 echoed this limitation: \enquote{I do not even know how much I have, I have to switch over again}. P06 demonstrated the operational risk of bypassing the table: placing a course from the graph blindly pushed a semester past its nine-ECTS cap, a breach only discovered upon returning to the table view. This observation explains why the final commit step naturally reverted to the table: users cannot confidently assign courses without immediate visual feedback on the temporal consequences of their placement.

\paragraph{What Does Not Change, and What Cannot Be Attributed}
\label{sec:eval-accuracy}

Placement accuracy remained stable, placing a strict operational bound on the impact of the substitution. The share of placement attempts ending in a system rejection shifted from a median of 16.7\,\% in Scenario~A to 13.6\,\% in Scenario~B. However, this share decreased for only seven participants and increased for four, a split consistent with chance (sign test, $p = 0.55$). The absolute number of rejections per ten minutes rose from 1.2 to 1.8, but that rate rises with placement activity and carries the same order confound as the placement rate itself; the share, a ratio internal to each run, is what bounds the substitution. Ultimately, the graph changed \emph{where} candidates were found, but left \emph{how} they were committed untouched. This stability is structurally logical: as detailed in Section~\ref{sec:eval-friction}, the rework loop at the commit stage is driven by obscured term parity and per-lane load metrics, which are properties of the table view unaffected by the graph.

\label{sec:eval-rate}
Because graph availability and second exposure are collinear in this fixed sequence (Section~\ref{sec:eval-order}), the remaining differences between the runs are documented here but left unattributed. Placement speed increased: placements per ten minutes rose from a median of 7.8 in Scenario~A to 13.9 in Scenario~B, in ten of eleven participants, while sampled time in the planner fell from a median of 20.8 to 13.3 minutes, despite the second task requiring an additional semester. Compliance improved: the six runs that met every constraint of their brief are all Scenario~B runs, and no Scenario~A run met its brief (Section~\ref{sec:eval-effectiveness}). Whether a participant became faster and more compliant due to the graph, interface familiarity, curriculum comprehension, or task repetition cannot be disentangled by this design. The lower block of Table~\ref{tab:ab} collects these quantities; the upper block holds the two measures the fixed order leaves interpretable, the catalogue displacement of Section~\ref{sec:eval-substitution} and the rejection share above.

\begin{table}[htpb]
    \centering
    \small
    \setlength{\tabcolsep}{5pt}
    \begin{tabular}{@{}>{\raggedright\arraybackslash}p{5.3cm}rrl r@{}}
        \toprule
        Measure & A & B & Direction & Sign Test \\
        \midrule
        \multicolumn{5}{@{}l}{\emph{Interpretable under the fixed order}} \\
        \quad Catalogue-only occupancy      & 27.4\,\% & 4.5\,\%  & down in 10 of 11 & \\
        \quad Rejected share of placements  & 16.7\,\% & 13.6\,\% & down in 7 of 11  & $p = 0.55$ \\
        \midrule
        \multicolumn{5}{@{}l}{\emph{Confounded with scenario order}} \\
        \quad Placements per 10 minutes     & 7.8      & 13.9     & up in 10 of 11 & \\
        \quad Sampled minutes per run       & 20.8     & 13.3     & down & \\
        \quad Rejections per 10 minutes     & 1.2      & 1.8      & up & \\
        \quad Runs meeting every constraint & 0 of 11  & 6 of 11  & all six in B & \\
        \bottomrule
    \end{tabular}
    \caption{Within-subject measures across the eleven participants, split by whether the fixed condition order permits an interpretation. Columns A and B are per-participant medians, except the last row, which counts runs.}
    \label{tab:ab}
\end{table}

\subsection{Around the Loop: Staging, Recommending, Configuring}
\label{sec:eval-around}

Three things sit beside the planning loop without being stations of it. The Parking Stage is what lets the two views work on one plan; the Recommendation Panel accompanies the loop without driving it; and one input to the checklist that steers the loop, the per-semester ECTS thresholds, was set by the students themselves. Each is reported with what it contributed and what it cost.

\paragraph{The Parking Stage: Shortlisting Before Commitment}
\label{sec:eval-parking}

The Parking Stage is what allows two surfaces to operate on one plan without either having to do the other's work, and every participant used it to shortlist before committing (E-G06, 11/11). In the graph-on condition it is where courses discovered in the graph were held before assignment. P04 described the strategy in advance (\enquote{I will put everything into the parking zone and then go back to the table and assign each to the right semester}), and the recordings show it executed: P09 moved nine courses from graph nodes into the Parking Stage in a single episode, and P10 nine across two.

That the stage sits outside the plan's accounting is both its function and a source of error. Parked credits are excluded from the planned ECTS total and nothing on screen says so (E-P38, 10/11). P02 planned against a total of 102~ECTS while 24~ECTS sat parked; P10 and P12 each ended a run with 6~ECTS parked while the Dashboard reported a complete 180~ECTS plan and zero missing requirements. Nine participants asked, in one form or another, for the stage to be visible from inside the graph, where it is not (E-P37, 9/11).

\paragraph{The Recommendation Panel Accompanies the Loop}
\label{sec:eval-recs}
%% REV (the thing we state its on ef the most occupied surfaces, but its not on our process decriped at the beginning. We have to make clear somewhere why)
Although the Recommendation Panel was one of the most occupied surfaces in the tool, it did not integrate into the primary planning cycle. It remained open for 24\,\% of sampled frames in Scenario~A and 20\,\% in Scenario~B, a share comparable to the catalogue in Scenario~A (27.1\,\%) and the graph in Scenario~B (26.6\,\%). Ten of the eleven participants opened it at least once; triangulation between the qualitative codebook (E-G02, 10/11) and the visual frame logs confirms the single participant who never accessed it.

High frame occupancy measures visual presence, which does not automatically equate to active consultation. The entry and dwell metrics in Table~\ref{tab:occupancy} distinguish these two behaviors. Active planning stations are characterized by frequent entries and short dwell times, indicating users navigated there for a specific query and left upon finding an answer: the catalogue was entered 40 times at an 81~s mean dwell, and the graph 42 times at 62~s. In stark contrast, the Recommendation Panel exhibited infrequent entries but exceptionally long continuous dwells (23 entries at a 125~s mean in Scenario~A, and just eight entries at a 244~s mean in Scenario~B). This skew is concentrated: 71\,\% of the panel's total open time occurred during just nine prolonged episodes lasting over three minutes. The most extreme case was P01, who opened the panel once and left it open for an unbroken 15 minutes while working on the canvas beneath it. Because the panel was left open as a persistent background state rather than repeatedly navigated to for active queries, it is methodologically classified as an accompanying overlay rather than a structural station of the planning loop.

Usage patterns diverged significantly based on participants' prior curriculum familiarity. Ten participants directly incorporated recommendations into their plans (E-G02, 10/11), and seven explicitly acknowledged the system dynamically re-ranking suggestions in response to their typed interests (E-G14, 7/11). Conversely, seven participants browsed the recommendations but ultimately bypassed them in favor of the deterministic checklist and catalogue search (E-N05, 7/11). Because six of these seven had successfully utilized recommendations earlier in the session, this reversion indicates a deliberate strategic choice rather than a failure of feature discovery. In contrast, the two participants lacking the domain knowledge %% REV (Trough what? Early in the program?)
to independently evaluate modules relied heavily on the panel and the Dashboard as compensatory mechanisms for their missing curriculum expertise (E-N06, 2/11).  The Scenario~B data quantifies this polarization: five participants never opened the panel, whereas three kept it open for over half the task duration.

Three usability defects contextualise these interaction patterns. First, because the panel operates as a UI overlay, it physically occluded the underlying semester lanes it was intended to populate, causing structural visual interference for six participants (E-P11, 6/11). Second, a recommendation family returned nothing where it was most needed: the internships tab rendered a bare empty state in a task that required an internship semester (E-P42, 4/11), which P09 met while searching for exactly that. The tab is keyed to the career direction stated in the profile and returns nothing when that field is blank (Sections~\ref{sec:rec-types-evidence} and~\ref{sec:profile}), so participants who had not filled it in were shown an empty tab rather than a prompt for the input it depends on. Finally, although every recommendation card provides a textual justification, four of the nine participants for whom this could be evaluated were unable to identify which of their specific profile inputs had triggered the suggestion (E-P68, 4/9). This indicates a gap in transparency.

\paragraph{User-Configured Compliance Thresholds}
\label{sec:eval-config}

Most of what the rule-checker enforces is fixed by the curriculum: module requirements, ECTS totals and sequencing warnings are properties of the programme rather than settings. The per-semester load is the exception. It is checked against two values exposed in the profile, a maximum and a recommended figure, defaulting to 42 and 30~ECTS; the first-use flow shows where they are, and aligning them to the brief was part of the task as set. Neither scenario band is expressible as a band, and a profile left untouched is checked against the defaults rather than against the brief, with nothing on screen to say which of the two is happening.

The recordings show both outcomes. P01 made the change within four minutes of first contact and discovered the asymmetry in the same breath (\enquote{I will set it to 33 straight away, more is not allowed, and is there a minimum? No}). P10 set 33 and 30 within three minutes. P12 narrated the mismatch in Scenario~B before correcting it by hand. P11 changed the recommended value from 30 to 24 midway through a run, which moved what the tool counted as a violation partway through the session. Where the profile was not aligned, the tool checked the wrong thing: P06 completed Scenario~B with one semester at 30~ECTS against a stated band of 21 to 27 while the Dashboard reported no warning, because the recommended value was still at its default of 30 and the warning fires only above it (E-P32, 11/11; E-P57, 11/11). Students trusted the verdict in both cases.

\subsection{What Worked and Where the Loop Breaks}
\label{sec:eval-findings}

Two structural patterns organize the evaluation findings. First, positive outcomes were universal: the five successful behavioral patterns evidenced across all eleven participants align directly with the core planning cycle detailed in Section~\ref{sec:eval-process}. Second, system failures were highly localized, clustering at three specific stages of this same cycle. An additional architectural defect operates independently of these stages: parked credits are excluded from the plan's overall totals without explicit system feedback (E-P38, 10/11). This section reports the universal positive findings, analyzes the friction points across the three workflow stages, and concludes with a prioritized usability problem register.

\paragraph{What Worked}
\label{sec:eval-positives}

Five positive findings are universal: the pre-filled plan removes the blank page (E-G01, 11/11), the requirement checklist drives completion (E-G15, 11/11), the Parking Stage holds candidates before commitment (E-G06, 11/11), real-time compliance feedback is noticed and used (E-G04, 11/11), and the graph's filters are valued once found (E-G07, 11/11). Table~\ref{tab:positives} lists these together with the less frequent findings, including the two recommendation codes analysed in Section~\ref{sec:eval-recs}.

\begin{table}[htpb]
    \centering
    \small
    \begin{tabular}{p{8.6cm}cp{2.4cm}}
        \toprule
        Positive finding (code) & Freq. & ISO dimension \\
        \midrule
        Pre-filled starting plan removes the blank page (E-G01) & 11/11 & efficiency/sat. \\
        Requirement checklist drives completion (E-G15) & 11/11 & effectiveness \\
        Parking Stage used to shortlist before committing (E-G06) & 11/11 & satisfaction \\
        Real-time compliance nudges noticed and used (E-G04) & 11/11 & efficiency \\
        Graph and sidebar filters valued (E-G07) & 11/11 & efficiency/sat. \\
        Recommendations actively used to fill gaps (E-G02) & 10/11 & efficiency/sat. \\
        Graph supports interest-driven browsing (E-G10) & 9/11 & efficiency/sat. \\
        Semester Picker offers only legal semesters (E-G03) & 7/11 & efficiency \\
        Recommendations respond to stated interests (E-G14) & 7/11 & satisfaction \\
        Graph clarifies module composition (E-G17) & 6/11 & satisfaction \\
        Graph \enquote{mandatory} filter shows affiliations (E-G16) & 4/11 & effectiveness \\
        \enquote{Done} check-off mode with live totals valued (E-G11) & 3/11 & satisfaction \\
        Visual legend aids interpretation (E-G13) & 3/11 & satisfaction \\
        Reordering plan sections brings structure (E-G12) & 3/11 & satisfaction \\
        Milestone animation creates delight (E-G05) & 2/11 & satisfaction \\
        \bottomrule
    \end{tabular}
    \caption{Positive findings across the eleven analysed participants.}
    \label{tab:positives}
\end{table}

\paragraph{Where the Loop Breaks}
\label{sec:eval-friction}

The most costly problems cluster at three stages of the loop, and each stage fails in a different way: the %% REV (is this introduces "drive stage")
drive stage does not show the status the loop depends on, the find stage cannot be narrowed to the gap being closed, and the commit stage withholds the fact needed to place a course correctly.

At the drive stage the status the loop depends on was not glanceable. Requirement and ECTS-gap status lived behind collapsible Dashboard panels (E-P10, 11/11), and where status was shown it was written as text rather than given as a figure (E-P03, 7/11; P01: \enquote{how many do I have? I do not see it}). A requirement named in the Dashboard's missing list could not be followed to the courses that satisfy it, so the module had to be re-found by hand in the catalogue (E-P58); all four of the most recent participants demonstrate the round trip, reading a module name out of the panel and retyping it into search. The warning list is the extreme case: across the 437 frames of the four most recent sessions the warnings panel was never opened once, while its counter climbed to between two and five in every run.

At the find stage the catalogue could not be narrowed. Filtering by category, obligation, or ECTS was available only in the graph's sidebar (E-P04, 11/11), so students in the graph-off condition fell back on free-text search. P10 states the fallback and its cost in one line: \enquote{that does not work, so I have to pick them out by hand}. The Recommendation Panel was the other route to a candidate, and Section~\ref{sec:eval-recs} reports both how far it carried students and where it returned nothing.

At the commit stage, the interface omits the temporal data required for accurate course placement. Because winter and summer availability is not rendered on the semester lanes, students were forced to guess the term, triggering predictable system rejections and immediate rework (E-P29, 11/11). This creates the study's most severe friction loop. Crucially, this is a visualization deficit rather than a validation error: the system evaluates the rule correctly, but withholds the information needed to prevent the error at the point of decision. To compensate, ten participants explicitly requested a term indicator on the lanes, while others resorted to calculating semester parity aloud.

\paragraph{The Problem Register}
\label{sec:eval-problems}

Table~\ref{tab:problems} ranks the highest-impact and most frequent problems across the eleven participants. Most belong to a stage of the loop and were developed above; the last column marks the five problems that have since been fixed. Those revisions were made after this study and are described in Chapter~\ref{chap:implementation} (Section~\ref{sec:impl-post-study}), so every problem the table lists was present in the system these participants used.

\begin{table}[htpb]
    \centering
    \small
    \begin{tabular}{@{}l>{\raggedright\arraybackslash}p{7.0cm}ccl c@{}}
        \toprule
        Code & Problem & Sev. & Freq. & ISO dimension & Fixed \\
        \midrule
        E-P04 & Table and catalogue lack the graph's filters & 3 & 11/11 & efficiency & \\
        E-P05 & Placement feedback arrives only after the drop & 2 & 11/11 & effectiveness & \\
        E-P10 & Requirement status behind collapsible panels & 2 & 11/11 & efficiency & \\
        E-P29 & Winter/summer offering not glanceable & 2 & 11/11 & efficiency & $\checkmark$ \\
        E-P32 & Workload warning driven by the profile, not the task band & 2 & 11/11 & efficiency & \\
        E-P25 & Reduced-load semester has no affordance & 2 & 11/11 & effectiveness & \\
        E-P26 & Workload band not enforced & 3 & 11/11 & effectiveness & \\
        E-P34 & Completion signalled over an incomplete constraint set & 3 & 11/11 & effectiveness & \\
        E-P53 & Placement menus omit and invent semesters & 2 & 11/11 & effectiveness & \\
        E-P57 & Workload target ignores the task band & 3 & 11/11 & efficiency & \\
        E-P38 & Parking silently removed from totals and counts & 2 & 10/11 & effectiveness & \\
        E-P35 & Lower workload bound only advisory & 3 & 10/11 & effectiveness & \\
        E-P27 & Study wording leaks into the profile UI & 1 & 10/11 & satisfaction & $\checkmark$ \\
        E-P47 & Graph carries no semester dimension & 3 & 10/10 & effectiveness & \\
        E-P48 & Module ECTS selector rewrites the plan total & 3 & 10/10 & effectiveness & \\
        E-P49 & Dashboard contradicts itself on a focus area's completion & 3 & 8/8 & effectiveness & \\
        E-P17 & Adding a semester not discoverable & 2 & 9/11 & effectiveness & \\
        E-P14 & Onboarding tour drives the early session & 1 & 9/11 & efficiency & \\
        E-P37 & Parking Stage not visible inside the graph & 2 & 9/11 & efficiency & \\
        E-P50 & Toast semantics unreliable & 2 & 9/9 & efficiency & $\checkmark$ \\
        E-P11 & Recommendation Panel occludes the semester lanes & 1 & 6/11 & efficiency & \\
        E-P02 & Totals and remaining ECTS not glanceable & 2 & 8/11 & efficiency & \\
        E-P03 & Requirement status stated only in prose & 2 & 7/11 & efficiency & \\
        E-P16 & Retuning the cap destabilises the plan & 3 & 7/11 & effectiveness & \\
        E-P01 & Curriculum and elective terminology unclear & 3 & 4/11 & effectiveness & \\
        E-P42 & Recommendation tabs render a bare empty state & 1 & 4/11 & satisfaction & $\checkmark$ \\
        E-P44 & Recommended ordering not shown at the point of placement & 2 & 4/11 & effectiveness & $\checkmark$ \\
        E-P68 & Recommendation provenance is opaque & 2 & 4/9 & satisfaction & \\
        \bottomrule
    \end{tabular}
    \caption{Highest-impact usability problems across the eleven analysed participants. Denominators below eleven mark codes that could not be decided for every participant (Section~\ref{app:eval-matrix}). A tick marks a problem addressed after the study (Section~\ref{sec:impl-post-study}).}
    \label{tab:problems}
\end{table}


\subsection{Questionnaire Results}
\label{sec:eval-questionnaire}

The User Experience Questionnaire (UEQ) results describe a tool participants judged capable before they judged it engaging: pragmatic quality ($+2.28$) outscores hedonic quality ($+2.00$), and attractiveness ($+2.61$) and efficiency ($+2.59$) are the strongest of the six scales (Table~\ref{tab:ueq}). The two weakest map onto the qualitative record. Perspicuity ($+1.82$) is the questionnaire's reflection of the learnability and terminology friction detailed in Section~\ref{sec:eval-friction}. Novelty ($+1.77$) reflects that participants read the tool as a better version of what they already did rather than as an unfamiliar concept; P04, who had planned in a spreadsheet, put it as \enquote{until now I only edited Excel files manually, I find this really great}. Dispersion sits in the same region: novelty (SD~0.78) and dependability (0.75) vary most across the cohort, attractiveness (0.33) and efficiency (0.30) least, so the cohort agreed on what the tool is for and differed on how far to trust it and how new it is. Per participant the overall score spans $+1.35$ (P06) to $+2.69$ (P05).

\begin{table}[htpb]
    \centering
    \small
    \begin{tabular}{lrrr}
        \toprule
        UEQ scale & Mean & SD & Range \\
        \midrule
        Attractiveness & $+2.61$ & 0.33 & $+2.17$ to $+3.00$ \\
        Efficiency     & $+2.59$ & 0.30 & $+2.00$ to $+3.00$ \\
        Dependability  & $+2.43$ & 0.75 & $+0.75$ to $+3.00$ \\
        Stimulation    & $+2.23$ & 0.68 & $+0.75$ to $+3.00$ \\
        Perspicuity    & $+1.82$ & 0.70 & $+0.75$ to $+3.00$ \\
        Novelty        & $+1.77$ & 0.78 & $+0.75$ to $+2.75$ \\
        \midrule
        Pragmatic quality & $+2.28$ & 0.45 & $+1.17$ to $+2.75$ \\
        Hedonic quality   & $+2.00$ & 0.64 & $+0.75$ to $+2.75$ \\
        Overall           & $+2.27$ & 0.39 & $+1.35$ to $+2.69$ \\
        \bottomrule
    \end{tabular}
    \caption{UEQ group profile across the eleven analysed participants ($n = 11$), on the standard $-3$ to $+3$ scale. Pragmatic quality averages Perspicuity, Efficiency and Dependability, hedonic quality averages Stimulation and Novelty, and Overall is the mean of all 26 items; Attractiveness enters neither aggregate. Values are descriptive; a sample of eleven does not support inferential claims.}
    \label{tab:ueq}
\end{table}



While the UEQ captured the system's global usability footprint, the Technology Acceptance Model (TAM) items administered after each task evaluate perceived utility across the two distinct conditions. Overall acceptance ratings remained consistently high across both scenarios (Table~\ref{tab:tam}). The quantitative differences between conditions are small. While both usefulness and ease of use increased in Scenario~B, the fixed-order design precludes attributing this strictly to the graph, as the shift remains collinear with second-task familiarity. In this exploratory context, the primary value of the TAM data lies in how it aligns with the behavioral observations regarding cohort bifurcation. Specifically, average satisfaction with the final plan decreased in Scenario~B (6.27 to 5.91). This decline does not appear uniform across the cohort; rather, it is driven primarily by the two participants for whom the graph introduced cognitive friction. Because these are the same individuals identified in the qualitative findings as lacking prerequisite curriculum knowledge.

\begin{table}[htpb]
    \centering
    \small
    \begin{tabular}{lrr}
        \toprule
        Construct (1--7) & Scenario~A & Scenario~B \\
        \midrule
        Perceived usefulness (3 items)        & 6.18 & 6.45 \\
        Perceived ease of use (2 items)       & 5.59 & 6.09 \\
        \midrule
        \enquote{Easy to keep an overview} (separate) & 6.36 & 6.55 \\
        Satisfaction with the plan             & 6.27 & 5.91 \\
        Task difficulty (7 = very easy)        & 6.18 & 6.00 \\
        \bottomrule
    \end{tabular}
    \caption{TAM-derived acceptance items and post-task ratings, displaying means across the eleven analysed participants. The \enquote{overview} item is reported separately rather than averaged into ease-of-use because it correlated more strongly with usefulness, which would compromise construct validity.}
    \label{tab:tam}
\end{table}

\subsection{The Plans Produced, Against Self-Report}
\label{sec:eval-effectiveness}

Every participant reported success in a condition in which no plan met its brief. All eleven rated their Scenario~A outcome as \enquote{Yes}, and not one Scenario~A run satisfied the stated constraints. Self-reported success in this study therefore measures felt rather than actual completion (E-N03, 9/11), and the observational record rather than the questionnaire carries the claim about what was produced. This is the reason the completion audit was conducted at all (Section~\ref{sec:eval-nongoals}), and it is the section's principal result.

The brief's constraints divide into two kinds, and the plans have to be read against each separately. The tool can represent the ECTS total, the curriculum's own requirements and a per-semester ceiling, the last through the profile a student sets; it cannot represent a lower bound on a semester's load or a reduced-load semester for an internship or exchange, because it has no minimum and no semester type. Against the constraints it can represent, seventeen of the twenty-one resolved runs are clean; the four that are not are two runs that ended short of the total (P01) and two with semesters above the ceiling planned against a profile still at its default (P03-B, P06-B), which is the configuration failure of Section~\ref{sec:eval-config}. Against the full brief, six of the twenty-two runs met every stated constraint, most of the remainder failed on one or two semesters, and one run is left unresolved because two lane values were never legible; the per-run audit is in Section~\ref{app:eval-completion}. The difference between seventeen and six is the two constraints the tool cannot hold: the reduced-load semester and the band's floor were handled by the students themselves, with mixed results and with no confirmation from the tool either way (E-P25 and E-P26, both 11/11). All six fully compliant runs are Scenario~B runs and Scenario~A produced none, which is the strongest single pattern in the completion data and is entirely confounded with order (Section~\ref{sec:eval-rate}).

The tool's own completion signal ran ahead of the observed state in the same way the self-reports did. All twenty-two runs ended with the Dashboard reporting zero missing requirements and a completed focus-area checklist, including a run with four lanes outside the stated band and no internship semester at all, and two runs with credits left in the Parking Stage. The calibration gap did not track experience: some participants over-rated, one was well calibrated, and an experienced participant under-rated a near-complete plan, so we report no consistent relationship between experience and calibration in this sample.

What the checklist did at the boundary of its own model is the sharpest instance of the same mechanism. P02 created an empty semester lane specifically to serve as the internship semester, then filled it to twelve ECTS because the Dashboard reported that six ECTS were still missing from the total and said nothing about the internship. The requirement he had represented himself was invisible to the mechanism he was following, so following the mechanism removed it. This is the study's central design implication, developed in Chapter~\ref{chap:discussion}: a checklist that students trust as their compass must cover every constraint they are held to, because the constraints it omits are the ones that will be optimised away.

\subsection{Summary}
\label{sec:eval-summary}
Eleven students planned a Bachelor curriculum twice, once without the curriculum graph and once with it. The two records support five conclusions about the planning process and one about how the tool was received. Each is taken up in Chapter~\ref{chap:discussion}.

Planning under this tool is a single checklist-driven loop. Students read what is missing, find a candidate, hold it, place it, and re-read. The loop was the completion mechanism for every participant, its two-state cycle accounts for roughly 40\,\% of all observed state changes in both conditions, and its period is the same with and without the graph.

The graph enters that loop at one point. It replaces the catalogue as the station at which a named gap is turned into candidate courses, taking almost exactly the occupancy and entry count the catalogue vacates, and it does not alter the rate at which placements are rejected. Its contribution is to the query step and not to the commit step, which is the division the design assigned in advance, and the reason it holds is structural: the graph renders no semester dimension, so the consequence of a placement made there is not visible from there.

What the graph makes legible is curriculum structure, and only the thinnest slice of sequence. Containment edges are drawn and valued, and the two prerequisite pairs the Bachelor curriculum fixes are drawn with them; the Compliance Engine reports those same two pairs as advisories when they are planned out of order, and one such advisory fired across the twenty-two runs. No other ordering information existed anywhere in the evaluated tool, in a view or in the engine. Two participants read the graph as making dependencies legible (E-G08), while four asked for ordering information at the point of placement and did not obtain it (E-P44). This evaluation therefore bounds its claim to the contribution of a browsable structural view carrying a near-empty prerequisite overlay, and cannot speak to what a view of recommended sequence would contribute, because none was on screen. This is where the second assumption of Section~\ref{sec:overall-assumptions}, that prerequisite sequencing is the planning-relevant structure a curriculum encodes, meets curricula that publish almost none; Chapter~\ref{chap:discussion} (Section~\ref{sec:disc-related}) draws the consequence. On the first assumption, the hybrid concept itself, the findings are narrower: participants used the two surfaces for different steps and seven of eleven articulated that division as the value of having both (E-N01, 7/11), but no non-hybrid alternative was tested, so whether one surface could have served both steps remains open.

Recommendations accompanied the loop without becoming part of it. The panel was open for about a fifth to a quarter of all sampled frames and ten of eleven participants used it, but it was entered a fraction as often as the catalogue or the graph, and seven participants browsed it and then returned to the checklist and search to decide. Its value divided by curriculum confidence: the participants who could not yet judge modules leaned on it, and those who could set it aside.

The tool was received well, and the questionnaires say where its weaknesses are rather than contradicting the process account. Every UEQ scale is positive and pragmatic quality (+2.28) leads hedonic (+2.00), so students experienced a capable instrument before an engaging one; the two weakest scales, perspicuity (+1.82) and novelty (+1.77), correspond to the terminology and learnability problems the observational record shows and to the fact that experienced participants judged the tool against spreadsheets they already knew. The acceptance items sit high in both conditions, and the small differences between conditions fall where the process account predicts without being attributable to the graph. What the questionnaires cannot carry is the question of what students produced, for the reason the next paragraph gives.

\TDmajor{This sentence depends on an unverified claim}{\enquote{The first of them have been fixed with regression tests} is the scoped version of a sentence that an earlier check could not support against the repository. Chapter~\ref{chap:implementation}, Section~\ref{sec:impl-post-study-fixes} carries the same claim and the ticked column of Table~\ref{tab:problems} depends on it. Verify all three against what is actually merged, and scope or delete together. \textbf{Checked on \texttt{main} of \texttt{hypridplanner}, 2026-09-06:} \texttt{docs/known-defects.md} lists the study's codes with none marked closed; no test under \texttt{frontend/tests} or \texttt{backend/tests} names an E-P code; the Recommendation Panel's empty state is the generic \enquote{No recommendations}; the course card does show a winter/summer marker from \texttt{termAvailability}; and the Bachelor Dashboard target is the literal 180 in \texttt{features/dashboard/metrics.ts}. Table~\ref{tab:problems} now ticks exactly the five codes Chapter~\ref{chap:implementation} names (the E-P11 tick had no counterpart there). Deferred with Chapter~\ref{chap:implementation}'s note, by your decision of 2026-09-05.}

The compliance layer certifies plans it has not checked. Its thresholds are supplied by the user rather than by the task, its lower bound only advises where its upper bound blocks, and two hard constraints of both briefs lie outside its model. All twenty-two runs ended with the tool reporting no missing requirements, and every participant reported success in a condition in which no run met the stated constraints. The instrument students trusted as their measure of completion was, for the constraints it does not model, silent rather than wrong.
The defects this study surfaced are recorded as a register in the tool's repository, ranked by the severity and frequency reported here, and the first of them have been fixed with regression tests that reproduce the observed failure; the remainder are open.
